\section{Week 5 -- Monte Carlo Simulation}
Monte Carlo simulation is used to simulate the behaviour of large systems by
sampling from the underlying probability distributions. It is a powerful tool
for solving complex problems that are difficult to solve analytically. In
finance, Monte Carlo simulation is commonly used for option pricing, risk
management, and portfolio optimization.

Especially useful where no closed-form solution exists, or when the problem is
too complex for analytical methods. It allows us to model and analyze systems
with a large number ofuncertainties and variables.

\textbf{Law of Large Numbers:} As the number of simulations increases, the average of
the simulated outcomes will converge to the expected value of the underlying
distribution.

$$\lim_{N \to \infty} \frac{1}{N} \sum_{i=1}^{N} X_i = \mathbb{E}[X]$$
The sample mean of the simulated outcomes will converge to the population mean as the number of simulations increases.

\textbf{Central Limit Theorem:} The distribution of the sample mean of a large number
of independent and identically distributed random variables will approximate a
normal distribution, regardless of the underlying distribution.

\textbf{Outcomes:} The results of the simulations can be used to estimate
probabilities, expected values, and other statistics of interest. The outcomes
can be visualized using histograms, scatter plots, or other types of charts to
gain insights into the behaviour of the system being simulated.

\begin{enumerate}
    \item \textbf{Define Domain:} Identify the problem you want to solve and define the domain of
          the simulation. This includes specifying the variables, parameters, and the
          range of possible values for each variable.
    \item \textbf{Generate Random Inputs:} Use random number generators to create random inputs
          for the simulation. These inputs should be drawn from the appropriate
          probability distributions that represent the uncertainties in the system.
    \item \textbf{Run Simulations:} Execute the simulation by running a large number of
          iterations. For each iteration, use the random inputs to calculate the outcome
          of the system based on the defined model or equations. This step is
          deterministic, as it follows a specific set of rules or equations to produce an
          outcome. The randomness comes from the random inputs generated in the previous
          step.
    \item \textbf{Analyze Results:} Use statistical methods to analyze the results of the
          simulations. This may include calculating means, variances, confidence
          intervals, or other relevant statistics. The outcomes can be visualized using
          histograms, scatter plots, or other types of charts to gain insights into the
          behaviour of the system being simulated.
\end{enumerate}

The accuracy of Monte Carlo simulation depends on the number of iterations and
the quality of the random inputs. More iterations generally lead to more
accurate results, but also require more computational resources. It is
important to balance the trade-off between accuracy and computational
efficiency when using Monte Carlo simulation.

Error decreases with the square root of the number of simulations, meaning that
to reduce the error by half, you need to quadruple the number of simulations.
This is known as the \textbf{square root law} of Monte Carlo simulation. Therefore, while increasing the
number of simulations can improve accuracy, it may not always be practical due
to computational constraints. In such cases, variance reduction techniques can
be employed to improve the efficiency of the simulation without necessarily
increasing the number of iterations.

$$\text{Error} \propto \frac{1}{\sqrt{N}}$$

The 95\% confidence interval for the estimated mean can be calculated using the
formula: $$\text{95\% CI} = \bar{X} \pm 1.96 \times \frac{\sigma}{\sqrt{N}}$$
where $\bar{X}$ is the sample mean, $\sigma$ is the standard deviation of the
sample, $N$ is the number of simulations, and $\frac{\sigma}{\sqrt{N}}$ is the
standard error. This interval provides a range within which we can be 95\%
confident that the true population mean lies.

\textbf{Monte Carlo Failure:} Monte Carlo simulation can fail to provide accurate
results if the necessary conditions are not met.

The necessary conditions for Monte Carlo simulation to work effectively
include:
\begin{enumerate}
    \item Proper distribution that matches the underlying process being simulated.
    \item Independece of simulations, ensuring that each simulation is independent of the
          others.
    \item Sufficient number of simulations to achieve the desired level of accuracy.
    \item Accurate model that correctly represents the system being simulated. If the
          formula or model used in the simulation is incorrect, the results will be
          unreliable regardless of the number of simulations.
\end{enumerate}

Note: The experience of any single path the simulation has taken can vary
wildy, and may not be representative of the overall distribution. Therefore, it
is important to run a large number of simulations to ensure that the results
are reliable and representative of the underlying distribution.

\subsection{Finance uses}
\textbf{Asian Options:} Monte Carlo simulation can be used to price Asian options,
which are options where the payoff depends on the average price of the
underlying asset over a certain period of time. The simulation can be used to
generate multiple price paths for the underlying asset, and the average price
can be calculated for each path to determine the payoff of the option. This
uses the Geometric Brownian Motion (GBM) model to simulate the price paths of
the underlying asset. The GBM model assumes that the price of the asset follows
a stochastic process with a constant drift and volatility. The formula for
simulating the price path of the asset is given by: $$S_{t+\Delta t} = S_t
    \exp\left((\mu - \frac{1}{2} \sigma^2) \Delta t + \sigma \sqrt{\Delta t}
    Z\right)$$ where $S_t$ is the price of the asset at time $t$, $\mu$ is the
drift, $\sigma$ is the volatility, $\Delta t$ is the time step, and $Z$ is a
standard normal random variable. $-\frac{1}{2} \sigma^2$ is the Ito correction
term that accounts for the volatility of the asset, ensuring that the expected
value of the simulated price paths matches the actual expected value of the
asset price.

\textbf{European Options:} While Black-Scholes model exists, we can also get the price
of an option via Monte Carlo simulation and check against Black-Scholes to
determine if our simulation is accurate and successful.

For a call option with strike price $K$, the payoff at maturity is given by:
$$\text{Payoff} = \max(S_T - K, 0)$$ where $S_T$ is the price of the underlying
asset at maturity. To price the option using Monte Carlo simulation, we can
simulate multiple price paths for the underlying asset using the GBM model,
calculate the payoff for each path, and then take the average of the payoffs to
estimate the option price. The estimated option price can be calculated using
the formula: $$\text{Option Price} = e^{-rT} \times \frac{1}{N} \sum_{i=1}^{N}
    \text{Payoff}_i$$ where $r$ is the risk-free interest rate, $T$ is the time to
maturity, $N$ is the number of simulations, and $\text{Payoff}_i$ is the payoff
for the $i$-th simulated price path. The estimated option price can then be
compared to the price obtained from the Black-Scholes model to assess the
accuracy of the Monte Carlo simulation. If the estimated price is close to the
Black-Scholes price, it suggests that the Monte Carlo simulation is successful
in capturing the dynamics of the underlying asset and providing a reasonable
estimate of the option price.